% concatenated from manuscript.tex
\documentclass{amsart}
\usepackage[margin=1in]{geometry}
\usepackage{amsmath,amssymb,amsthm}
\usepackage{hyperref}

\newtheorem{theorem}{Theorem}[section]
\newtheorem{proposition}[theorem]{Proposition}
\newtheorem{lemma}[theorem]{Lemma}
\newtheorem{corollary}[theorem]{Corollary}
\newtheorem{remark}[theorem]{Remark}

\newcommand{\T}{\mathbb T}

\title{Uniqueness of synchronized stationary equilibria in the Kuramoto mean field game}
\author{Sebastian Munoz}
\address{Department of Mathematics, University of California, Los Angeles, Los Angeles, CA 90095, USA}
\email{sebastian@math.ucla.edu}
\subjclass[2020]{91A16, 35Q89, 49N80, 35B32, 35J60, 60H30}
\keywords{Kuramoto model, mean field games, synchronization, stationary equilibrium, bifurcation, Hamilton--Jacobi--Bellman equation, uniqueness}
\date{}

\begin{document}

\begin{abstract}
The stationary Kuramoto mean field game models a population of phase
oscillators that form synchronized Nash equilibria above a critical
interaction strength. We prove that the synchronized branch is a
unique smooth family of Nash equilibria emerging from the uniform
state at the bifurcation: at each supercritical interaction strength
the synchronized equilibrium is unique up to rotation of the torus,
and converges smoothly to the uniform distribution as
the interaction parameter decreases to the critical threshold. Both follow
from our main technical result: the scalar self-consistency map is strictly concave, settling a
conjecture of~\cite{CCS}. The proof decomposes the second derivative
of the self-consistency map into two sign-indefinite moments of the
equilibrium---a cubic moment and a gradient moment---and controls
their signs through sharp shape estimates for the value function, a
pointwise geometric-mean monotonicity that determines the sign of the
cubic moment via a cosine-skewness inequality, and a reflection argument combined with
a correlation inequality for the gradient moment.
\end{abstract}

\maketitle

\section{Introduction}\label{sec:intro}

In a non-monotone mean field game, multiple nontrivial Nash equilibria
can a priori coexist, and a central structural question is whether a
given phase transition selects a unique branch organized by the model
parameters. We resolve this question for
the Kuramoto mean field game studied by Carmona, Cormier, and
Soner~\cite{CCS}, a mean field game
model of phase synchronization. The underlying Kuramoto
model~\cite{Kuramoto,Strogatz,ABPVS} is a paradigm
of phase-transition synchronization: a population of coupled phase
oscillators that remains incoherent below a critical coupling strength
and self-organizes into a synchronized state above that threshold. In the classical
mean-field limit, where the population evolves according to a
prescribed McKean--Vlasov dynamics, the phase transition and the
stability of the incoherent state are by now well understood at the
level of both the nonlinear Fokker--Planck equation~\cite{GPP} and the
fluctuations of the underlying interacting particle system~\cite{BGP}.

A mean field game~\cite{LL,HMC,CarmonaDelarue} version of the model,
in which each oscillator chooses its dynamics to minimize a cost
rather than following a prescribed law of motion, was studied
formally in~\cite{YMMS} and used in~\cite{CG} for an analysis of
jet-lag recovery. Related bifurcation-theoretic and global techniques for
non-monotone mean field games appear in~\cite{Cirant}. The first
rigorous analysis of the resulting stationary mean field game was
given in~\cite{CCS}, where it was shown that the
synchronization phase transition survives in the game-theoretic
setting as a bifurcation of the Nash equilibria, with full
synchronization in the strong-coupling limit $\kappa\to\infty$. This
framework has since been extended to heterogeneous populations with
random intrinsic frequencies~\cite{CCS2}, and adapted in~\cite{HS} to
a finite-state synchronization mean field game exhibiting the same
phase-transition phenomenology.

In the model we study, the representative oscillator is a controlled
diffusion on the torus $\T:=\mathbb R/2\pi\mathbb Z$,
\[
dX_t=\alpha_t\,dt+\sigma\,dB_t,
\]
in which the agent chooses the drift $\alpha$ to minimize the
discounted cost
\[
J(\alpha):=\mathbb E\!\left[\int_0^\infty e^{-\beta t}
\Bigl(\kappa\,c(X_t,\mu_t)+\tfrac12\alpha_t^2\Bigr)dt\right]
\]
against the prevailing population distribution
$\mu_t\in\mathcal P(\T)$. Here $\beta>0$ is the discount rate,
$\sigma>0$ the noise level, $\kappa>0$ the interaction strength, and
\[
c(\theta,\mu):=\int_\T 2\sin^2\!\tfrac{\theta-\theta'}{2}\,\mu(d\theta')
=1-\mu\bigl(\cos(\theta-\cdot)\bigr)
\]
is the classical Kuramoto coupling, which penalizes misalignment of the
agent's phase with the population. For a prescribed time-dependent
population flow $(\mu_t)_{t\geq 0}$, let $X^{t,x,\alpha}$ denote the
controlled process started from $x$ at time $t$, and define the value
function
\[
\Phi(t,x):=\inf_\alpha\mathbb E\!\left[\int_t^\infty e^{-\beta(s-t)}
\Bigl(\kappa\,c(X_s^{t,x,\alpha},\mu_s)
+\tfrac12\alpha_s^2\Bigr)ds\right],
\]
where $dX_s^{t,x,\alpha}=\alpha_s\,ds+\sigma\,dB_s$ and
$X_t^{t,x,\alpha}=x$. Dynamic programming gives the feedback
$\alpha_s^*=-\Phi_x(s,X_s)$ and the Hamilton--Jacobi--Bellman (HJB) equation
\[
-\Phi_t+\beta\Phi-\frac{\sigma^2}{2}\Phi_{xx}
+\tfrac12\Phi_x^2=\kappa\,c(x,\mu_t).
\]

A \emph{Nash equilibrium} is a population flow $(\mu_t)_{t\geq 0}$ such
that, when the representative oscillator uses this optimal feedback,
the law of the optimally controlled state is exactly the prescribed
population flow: $\mathcal L(X_t^*)=\mu_t$ for all $t\geq 0$. The full
time-dependent equilibrium problem is therefore a forward--backward
fixed point: the HJB equation couples $\Phi$ at time $t$ to the future
of $(\mu_s)_{s\geq t}$, while the controlled law
$\mathcal L(X_t^*)$ propagates forward in $t$.

In this paper we focus on stationary equilibria, of the form
$\mu_t\equiv\mu$. These are the game-theoretic analogue of the
incoherent and synchronized phases of the classical Kuramoto model:
they describe possible long-run organized states of the population,
and the phase transition observed in ~\cite{CCS} occurs through a bifurcation of
such stationary Nash equilibria. Uniqueness and continuity of the
stationary synchronized branch are thus basic structural questions,
prior to any analysis of stability or non-stationary selection.

In the stationary case the infinite-horizon control problem is
time-homogeneous, so $\Phi(t,x)=\phi(x)$, where
$\phi\in C^\infty(\T)$ solves the stationary
HJB equation
\[
\beta\phi-\frac{\sigma^2}{2}\phi_{xx}+\tfrac12\phi_x^2=\kappa\,c(\cdot,\mu),
\]
and the optimal feedback is $\alpha_t^*=-\phi_x(X_t)$. The stationary
equilibrium condition is that $\mu$ be invariant for the optimal
diffusion $dX_t=-\phi_x(X_t)\,dt+\sigma\,dB_t$, meaning that
$X_0\sim\mu$ implies $X_t\sim\mu$ for all $t\geq 0$. This invariant
law has density
\[
\mu(dx)\propto \exp(-2\phi(x)/\sigma^2)\,dx \qquad\text{on }\T.
\]
The interaction $c$ is rotation-invariant, so any rotation of an equilibrium is itself an
equilibrium. Rotating so that $\mu(\sin)=0$ and $\mu(\cos)\geq 0$
(with $\mu(\cos)>0$ in the synchronized case), and absorbing the constant
$\kappa$ on the right-hand side of the HJB into $\phi$, the
self-consistency condition reduces to
\begin{equation}\label{eq:self-consistency}
\mu=\mu^{\kappa\,\mu(\cos)},
\end{equation}
where, for each $\gamma\geq 0$, $\mu^\gamma$ is the probability measure
on $\T$ defined analogously from $\phi^\gamma\in C^\infty(\T)$, the
unique smooth solution of
\begin{equation}\label{eq:HJB-intro}
\beta\phi^\gamma-\frac{\sigma^2}{2}\phi^\gamma_{xx}
+\tfrac12(\phi^\gamma_x)^2=-\gamma\cos x,\qquad x\in\T.
\end{equation}
With this normalization,~\eqref{eq:self-consistency} is the fixed-point equation
$\gamma=F_\kappa(\gamma)$ on $[0,\infty)$ for the scalar
\emph{self-consistency map}
\[
F_\kappa(\gamma):=\kappa\,\mu^\gamma(\cos),
\]
whose value at the order parameter $\gamma=\kappa\mu(\cos)$ measures
the population-averaged force pulling each agent toward synchrony.
The choice $\gamma=0$ recovers the uniform measure, which is therefore
a stationary equilibrium for every $\kappa$. We refer to non-uniform
stationary equilibria as \emph{synchronized}. A short calculation,
carried out in Section~\ref{sec:setup}, gives
$F_\kappa(0)=0$, $F_\kappa'(0)=\kappa/\kappa_c$, and
$|F_\kappa|\leq\kappa$, where the \emph{critical interaction parameter}
is
\[
\kappa_c:=\beta\sigma^2+\frac{\sigma^4}{2}.
\]
In~\cite{CCS} it is shown that a synchronized stationary equilibrium
exists for every $\kappa>\kappa_c$, and that any sequence of such equilibria converges
up to rotation to $\delta_0$ (full synchronization) as
$\kappa\to\infty$. In the subcritical regime, global attraction of the uniform state is established
in the smaller weak-interaction regime $\kappa<\beta\sigma^2/4$, while for all
$\kappa<\kappa_c$, they construct equilibrium flows, starting from initial
laws sufficiently close to the uniform distribution, that converge exponentially
to the uniform state. In the numerical experiments of~\cite{CCS},
$F_\kappa$ was always concave and the synchronized equilibrium
was unique up to rotation. On this basis it was
conjectured~\cite[Rem.~7.4]{CCS} that $F_\kappa$ is concave on
$[0,\infty)$ for every $\beta,\sigma,\kappa>0$. Two structural
consequences sought there follow at once from this concavity: the
synchronized equilibrium would be unique up to rotation, and would
converge to the uniform measure at the bifurcation $\kappa=\kappa_c$.

The first partial result toward this conjecture was obtained by
Cesaroni and Cirant~\cite{CesaroniCirant}, who work with a closely
related model:
the ergodic Kuramoto mean field game, in which the discount
factor $\beta>0$ of~\eqref{eq:HJB-intro} is replaced by an ergodic
constant determined together with the value function and the
equilibrium measure. For that model it is shown that, up to rotation,
the synchronized stationary equilibrium is unique provided the
interaction parameter $\kappa$ is sufficiently large, together with a
local stability/turnpike result for even finite-horizon equilibria in
the same regime. Their argument is variational and uses
energetic estimates that are sharp in the strong-coupling limit.
It leaves open the interval between the corresponding bifurcation
threshold and their large-coupling threshold, and in particular gives
no information about the synchronized branch near the bifurcation,
which is the regime in which the phase transition itself occurs. The full concavity conjecture and
the behavior of the synchronization branch at the threshold have
remained open in either formulation. 


\begin{theorem}\label{thm:main}
For every $\beta,\sigma,\kappa>0$:
\begin{enumerate}
\item[\textup{(i)}] the map $F_\kappa$ is strictly concave on
$[0,\infty)$.
\item[\textup{(ii)}] for $\kappa\leq\kappa_c$ the stationary Kuramoto
mean field game admits no synchronized stationary equilibrium, while
for $\kappa>\kappa_c$ it admits, up to rotations of $\T$, a unique
synchronized stationary equilibrium $\mu_\kappa$ (normalized by
$\mu_\kappa(\sin)=0$ and $\mu_\kappa(\cos)>0$). Moreover, the density
of $\mu_\kappa$ depends $C^\infty$-smoothly on $\kappa\in(\kappa_c,\infty)$
as a map into $C^\infty(\T)$, and converges to the uniform density
$1/(2\pi)$ in $C^\infty(\T)$ as $\kappa\downarrow\kappa_c$.
\end{enumerate}
\end{theorem}

Theorem~\ref{thm:main}
resolves the conjecture of~\cite{CCS}, including in the regime near
the bifurcation that lies outside the reach of the energetic methods
of~\cite{CesaroniCirant}. 

Part~(i) is the main analytical result of the theorem: it proves the
concavity conjecture of~\cite[Rem.~7.4]{CCS}. Part~(ii) is a short
scalar consequence: the strictly concave function
$g_\kappa(\gamma):=F_\kappa(\gamma)-\gamma$ satisfies $g_\kappa(0)=0$,
$g_\kappa'(0)=\kappa/\kappa_c-1$, and $g_\kappa(\gamma)\to-\infty$ as
$\gamma\to\infty$, so it has exactly one positive zero
$\gamma_*(\kappa)$ for $\kappa>\kappa_c$, and this zero tends to $0$
as $\kappa\downarrow\kappa_c$. The deduction is carried out at the
end of Section~\ref{sec:proof}.

We expect the same method, with minor modifications, to
resolve the analogous uniqueness question for the ergodic Kuramoto
mean field game of~\cite{CesaroniCirant}. In their parametrization the
analog of part~(i) would be strict \emph{convexity} of the scalar
self-consistency map throughout the post-bifurcation regime, which
would give uniqueness, up to rotation, of the synchronized stationary
equilibrium for every supercritical interaction strength. The argument
should in fact be shorter in that setting, since the delicate cubic
moment discussed below is absent in the ergodic problem.

The remainder of the paper is devoted to the proof of part~(i).
Passing to the dimensionless variables of Section~\ref{sec:setup}, in
which $\zeta$ rescales $\gamma$ and $v=v^\zeta$ rescales $\phi^\gamma$,
reduces the claim to strict concavity of $A(\zeta):=\mu^\zeta(\cos)$ in
$\zeta>0$, where $\mu^\zeta(dx)=f(x)\,dx$ has density $f\propto e^{-v}$
and $v\in C^\infty(\T)$ solves
\[
-v_{xx}+\lambda v+\tfrac12 v_x^2=-\zeta\cos x,
\qquad\lambda:=2\beta/\sigma^2.
\]
Writing
\[
w:=A/\lambda+v_\zeta,\qquad z:=v_{x\zeta},
\]
Lemma~\ref{lem:second-deriv} expresses the second derivative as
\begin{equation}\label{eq:second-derivative-intro}
A''(\zeta)=-\lambda\int_\T w^3 f\,dx
-3\int_\T wz^2 f\,dx,\qquad \int_\T wf\,dx=0,
\end{equation}
the additive constant $A/\lambda$ in $w$ chosen so that $w$ is
$\mu^\zeta$-centered. Equivalently, $-w=\partial_\zeta\log f^\zeta$ is
the score function of the Gibbs family $\{\mu^\zeta\}$, and the
centering is its standard zero-mean property. Both terms on the right
are then cumulants of $w$ under $\mu^\zeta$: the \emph{cubic moment}
$\int_\T w^3 f\,dx$ (the third cumulant of $w$) and the \emph{gradient
moment} $\int_\T wz^2 f\,dx$ (the covariance of $w$ with $z^2$).
Neither integrand has an obvious sign, and direct integration-by-parts
estimates do not give a sign for either.

Both moments are controlled instead through sharp shape estimates for
$v$ and $v_\zeta$ along $(0,\pi)$. On that interval,
$u:=v_x$ is a positive concave bump that leans toward $\pi$:
in addition to concavity, the quotient $u/\sin x$ is monotone
nondecreasing (Lemma~\ref{lem:u-shape}). The $\zeta$-derivative obeys a
related shape estimate after rescaling: setting
\[
Z:=\frac{z}{\sin x},
\]
the quotient satisfies $0\leq Z_x\leq uZ$ on $(0,\pi)$
(Lemma~\ref{lem:Z-bounds}). The natural
target is concavity of $u$, but a maximum-principle attack on $u_{xx}$
fails: the equation it satisfies has a zero-order term of uncontrolled
sign. The key step in Lemma~\ref{lem:u-shape} is to pair
concavity with the less obvious monotonicity of $u/\sin x$ and prove
both simultaneously with a bootstrapping argument via an auxiliary coupled ODE system.
The two quantities tracked are the Wronskian $\mathcal W$ of $u$ against the
forcing mode $\sin x$ (encoding monotonicity of $u/\sin x$) and the
normalized curvature $\mathcal K$ defined by $u_{xx}=-u\mathcal K$ (encoding concavity
of $u$). Because the forcing is the first Dirichlet mode, $\mathcal W$ and
$\mathcal K$ close into a first-order system.

These shape estimates reduce each moment
in~\eqref{eq:second-derivative-intro} to a one-dimensional inequality,
by substantially different arguments. The non-negativity of the cubic moment is the more difficult estimate and requires further insight beyond the shape information. An algebraic substitution rewrites the integrand as
a cubic in a single variable: writing $a:=v_\zeta$, $\Delta_0(x):=a(x)-a(0)$,
$\Delta_\pi(x):=a(\pi)-a(x)$, and $\Delta:=a(\pi)-a(0)$, the centered quantity
$A/\lambda+a$ is affine in $C:=(\Delta_\pi-\Delta_0)/\Delta$, and setting
$C=\cos\theta$ pushes $f$ forward to a density $\eta$ on $[0,\pi]$
under which the cubic moment becomes the centered cubic
$\cos\theta$-moment of $\eta$. Two independent facts then give it a sign:
$\eta$ is nonincreasing (Proposition~\ref{prop:cubic}), and any
nonincreasing density on $[0,\pi]$ has nonnegative centered cubic
$\cos$-moment, by a symmetrization argument across the median
(Lemma~\ref{lem:monotone-cosine-skewness}). Monotonicity of $\eta$ is
the substantive step. It does not follow directly from the shape
estimates, and is equivalent to a pointwise inequality on
$(0,\pi)$ that mixes $u$, the curvature of $C$, and $C$ itself. We
deduce it from a strictly stronger pointwise statement that may be of
independent interest: the geometric mean
\[
H(x):=e^{-v(x)}\,\frac{\sqrt{\Delta_0(x)\,\Delta_\pi(x)}}{\sin x}
\]
of the one-sided increments of $a$, normalized by $\sin x$ and
weighted by the Gibbs factor, is itself nonincreasing on $(0,\pi)$
(Lemma~\ref{lem:geometric-mean}).
The proof of this geometric-mean monotonicity rests on sharp integral
comparisons for $\Delta_0,\Delta_\pi$ that simultaneously exploit the concavity and
lean of $u$ from Lemma~\ref{lem:u-shape} and the bounds on $Z$ from
Lemma~\ref{lem:Z-bounds}. For the gradient
moment, reflection across $x=\pi/2$ proves a central-positivity
statement for $w$ (Lemma~\ref{lem:central-positivity}): under the
equilibrium $f$, $w$ has positive integral against every test weight
 symmetric about $\pi/2$ and nondecreasing on $[0,\pi/2]$. With
$w$ and $Z^2$ nondecreasing on $(0,\pi)$, applying this central
positivity together with a correlation
inequality yields a sign on the full gradient moment
(Proposition~\ref{prop:gradient}).

\section{Setting and dimensionless reduction}\label{sec:setup}

To avoid carrying $\beta$ and $\sigma$ explicitly through the proof,
we work in dimensionless variables. Setting
\[
\lambda:=\frac{2\beta}{\sigma^2},\qquad
\zeta:=\frac{4\gamma}{\sigma^4},\qquad
v^\zeta:=\frac{2}{\sigma^2}\phi^\gamma ,
\]
the HJB equation~\eqref{eq:HJB-intro} becomes
\begin{equation}\label{eq:HJB}
-v^\zeta_{xx}+\lambda v^\zeta+\tfrac12(v^\zeta_x)^2=-\zeta\cos x,
\qquad x\in\T,
\end{equation}
and the law $\mu^\gamma$ has density proportional to $e^{-v^\zeta}$. We will denote this measure by $\mu^\zeta$ in the new variable (the precise normalization is given in~\eqref{eq:mu} below). Thus
\[
F_\kappa(\gamma)=\kappa\,A\!\left(\frac{4\gamma}{\sigma^4}\right),
\qquad A(\zeta):=\mu^\zeta(\cos).
\] Equivalently, in the $\zeta$ coordinate the
fixed-point equation $\gamma=F_\kappa(\gamma)$ becomes
\[
\zeta=\frac{4\kappa}{\sigma^4}A(\zeta).
\]
The map $\gamma\mapsto\zeta$ is a positive linear rescaling, so strict
concavity of the original self-consistency map $F_\kappa$ in $\gamma$
is equivalent to strict concavity of $A$ in $\zeta$.

For the rest of the paper we fix $\lambda>0$. For completeness, we briefly establish existence and uniqueness of smooth solutions to the HJB equation, together with smooth dependence on the parameter $\zeta$. We write $C^{k,\alpha}(\T)$ for the Banach space of $C^{k,\alpha}$ functions on the torus $\T$, equivalently $2\pi$-periodic $C^{k,\alpha}$ functions on $\mathbb R$, with the usual H\"older norm on one period.

\begin{proposition}[Well-posedness and smooth dependence]\label{prop:smooth-zeta}
For every $\zeta\geq 0$, equation~\eqref{eq:HJB} has a unique solution
$v^\zeta\in C^\infty(\T)$. This solution is even with respect to both
$0$ and $\pi$, and every partial derivative $\partial_\zeta^j\partial_x^k v^\zeta$ is
even with respect to $0$ and $\pi$ when $k$ is even, and odd with respect to $0$ and $\pi$ when $k$
is odd. Moreover, for every $k\geq 0$ and $\alpha\in(0,1)$, the branch
$\zeta\mapsto v^\zeta$ is $C^\infty$ as a map into
$C^{k,\alpha}(\T)$.
\end{proposition}
\begin{proof}
We first prove existence on bounded parameter intervals. Fix
$M>0$ and $\alpha\in(0,1)$, and let $\mathcal I_M\subset[0,M]$ be the set
of $\zeta$ for which~\eqref{eq:HJB} admits a solution in
$C^{2,\alpha}(\T)$. Since $v\equiv0$ solves~\eqref{eq:HJB} at
$\zeta=0$, the set $\mathcal I_M$ is nonempty.

We record the local continuation statement that will be used both for
openness and for smooth dependence. Suppose that $v$ solves
\eqref{eq:HJB} at $\zeta_0\in\mathcal I_M$, and define
\[
\mathcal F(v,\zeta):=-v_{xx}+\lambda v+\frac12 v_x^2+\zeta\cos x
\]
as a map from
$C^{2,\alpha}(\T)\times\mathbb R$ to
$C^{0,\alpha}(\T)$. The derivative in the $v$ variable is
\[
D_v\mathcal F(v,\zeta_0)h=-h_{xx}+v_xh_x+\lambda h.
\]
The maximum principle gives a trivial kernel, since $\lambda>0$. This
operator is Fredholm of index zero: if $L_0=-\partial_{xx}+1$, then
\[
D_v\mathcal F(v,\zeta_0)L_0^{-1}=I+K,\qquad
K=(v_x\partial_x+\lambda-1)L_0^{-1},
\]
and $K$ is compact on $C^{0,\alpha}(\T)$ by the compact
embedding $C^{1,\alpha}(\T)\hookrightarrow
C^{0,\alpha}(\T)$. Thus $D_v\mathcal F(v,\zeta_0)$ is an
isomorphism by the Fredholm alternative~\cite[Thm.~5.3]{GT}. The
implicit function theorem gives neighborhoods $U$ of $v$ and $J$ of
$\zeta_0$, and a $C^\infty$ map
$\zeta\in J\mapsto \tilde v^\zeta\in U$, such that
$\mathcal F(\tilde v^\zeta,\zeta)=0$. In particular,
$J\cap[0,M]\subset\mathcal I_M$, so $\mathcal I_M$ is open.

We show that $\mathcal I_M$ is closed. Let
$\zeta_n\in\mathcal I_M$ with $\zeta_n\to\zeta\in[0,M]$, and let $v_n$
be corresponding solutions. The maximum principle applied to~\eqref{eq:HJB} gives
\[
\|v_n\|_{L^\infty}\leq \frac M\lambda .
\]
Writing $u_n:=v_{n,x}$ and differentiating~\eqref{eq:HJB}, we have
\[
-u_{n,xx}+\lambda u_n+u_nu_{n,x}=\zeta_n\sin x,
\]
and the maximum principle applied to $u_n$ gives
$\|u_n\|_{L^\infty}\leq M/\lambda$. Returning to~\eqref{eq:HJB},
\begin{equation}\label{eq:vnxx}
v_{n,xx}=\lambda v_n+\frac12 u_n^2+\zeta_n\cos x
\end{equation}
is uniformly bounded. Hence $(v_n)$ is precompact in $C^1(\T)$. Along
a subsequence, $v_n\to v$ and $v_{n,x}\to v_x$ uniformly. Equation~\eqref{eq:vnxx}
then gives uniform convergence of $v_{n,xx}$, and
the limit $v\in C^2(\T)$ solves~\eqref{eq:HJB} at the parameter
$\zeta$. Schauder estimates upgrade $v$ to $C^{2,\alpha}(\T)$,
so $\zeta\in\mathcal I_M$. Thus
$\mathcal I_M=[0,M]$. Since $M$ was arbitrary, solutions exist for every
$\zeta\geq0$.

Uniqueness follows from the maximum principle. If $v_1$ and $v_2$ solve
\eqref{eq:HJB} at the same $\zeta$, then $\psi:=v_1-v_2$ satisfies
\[
-\psi_{xx}+\frac12(v_{1,x}+v_{2,x})\psi_x+\lambda\psi=0,
\]
so $\psi\equiv0$. Since~\eqref{eq:HJB} is invariant under $x\mapsto-x$,
uniqueness gives $v^\zeta(-x)=v^\zeta(x)$. Combining this with
periodicity gives
$v^\zeta(\pi+y)=v^\zeta(-\pi-y)=v^\zeta(\pi-y)$, so $v^\zeta$ is also
even about $\pi$. The stated parities of the mixed partials follow by
differentiating the identities $v^\zeta(-x)=v^\zeta(x)$ and
$v^\zeta(\pi+y)=v^\zeta(\pi-y)$ in $\zeta$ and $x$.

It remains only to identify the regularity of the global branch. The
local continuation statement above applies at every parameter
$\zeta_0\geq0$, now with $v=v^{\zeta_0}$. By uniqueness, the local
branch $\tilde v^\zeta$ agrees with the globally defined solution
$v^\zeta$ wherever both are defined. Hence
$\zeta\mapsto v^\zeta$ is $C^\infty$ as a map into
$C^{2,\alpha}(\T)$. Bootstrapping upgrades this to smoothness
into each $C^{k,\alpha}(\T)$.
\end{proof}

The stationary equilibrium measure $\mu^\zeta\in\mathcal P(\T)$ is defined
by
\begin{equation}\label{eq:mu}
\mu^\zeta(dx)=f^\zeta(x)\,dx,
\qquad
f^\zeta(x):=\frac{e^{-v^\zeta(x)}}{\mathcal Z^\zeta},
\qquad
\mathcal Z^\zeta:=\int_\T e^{-v^\zeta(y)}\,dy.
\end{equation}
The first cosine moment of $\mu^\zeta$ defines
\begin{equation}\label{eq:A}
A(\zeta):=\mu^\zeta(\cos)=\int_\T \cos x\,f^\zeta(x)\,dx.
\end{equation}
We integrate over $\T$ with respect to $dx$ throughout and pass freely
between $[0,\pi]$ and $\T$ using evenness of $v^\zeta$ and $f^\zeta$ about
$0$.

A first-order computation in $\zeta$ identifies the bifurcation
threshold. At $\zeta=0$ the unique solution of~\eqref{eq:HJB} is
$v^0\equiv 0$, so by Proposition~\ref{prop:smooth-zeta} the function
$w:=\partial_\zeta v^\zeta|_{\zeta=0}\in C^\infty(\T)$ is well-defined.
Differentiating~\eqref{eq:HJB} in $\zeta$ at $\zeta=0$ gives
\[
-w_{xx}+\lambda w=-\cos x,
\qquad\text{so}\qquad w(x)=-\frac{\cos x}{\lambda+1}.
\]
Differentiating~\eqref{eq:A} at $\zeta=0$ with $f^0\equiv 1/(2\pi)$,
the contribution from $\partial_\zeta\mathcal Z^\zeta$ vanishes since
$\int_\T\cos x\,dx=0$, and
\[
A'(0)=-\int_\T \cos x\,w(x)\,\frac{dx}{2\pi}
=\frac{1}{2\pi(\lambda+1)}\int_\T\cos^2 x\,dx
=\frac{1}{2(\lambda+1)}.
\]
Translating back via $F_\kappa(\gamma)=\kappa A(4\gamma/\sigma^4)$ and
$\kappa_c=\sigma^4(\lambda+1)/2$,
\begin{equation}\label{eq:F-properties}
F_\kappa(0)=0,\qquad F_\kappa'(0)=\frac{\kappa}{\kappa_c},\qquad |F_\kappa|\leq\kappa,
\end{equation}
where the last bound follows from $|\cos|\leq 1$.

In dimensionless variables, Theorem~\ref{thm:main}(i) is the statement
that $A''(\zeta)<0$ for every $\lambda>0$ and every $\zeta>0$,
which we prove in Sections~\ref{sec:prelim}
and~\ref{sec:proof}.

\section{Shape estimates}\label{sec:prelim}

Throughout this section we fix $\lambda>0$ and $\zeta>0$ and write
$v=v^\zeta$, $u=v_x=v^\zeta_x$, $f=f^\zeta$,
$\mathcal Z=\mathcal Z^\zeta$.

The first lemma is the quantitative shape statement for $u$ on
$(0,\pi)$. Geometrically, $u$ is a positive concave bump that
leans toward $\pi$, and its sine-rescaling $u/\sin x$ is monotone. Beyond the four shape inequalities, the lemma also
records the monotonicity $\mathcal B_x\leq 0$ of the combination
$\mathcal B:=1+\lambda+u_x+u\cot x$. This $\mathcal B$ arises as the
zero-order coefficient of an elliptic equation for
$v_{x\zeta}/\sin x$ in Lemma~\ref{lem:Z-bounds}, and the monotonicity
$\mathcal B_x\leq 0$ is precisely what allows us to bound that
quotient there.
\begin{lemma}[Shape of $u$]\label{lem:u-shape}
On $(0,\pi)$,
\begin{equation}\label{eq:shape}
u>0,\qquad u_{xx}\leq 0,\qquad
\Bigl(\frac{u}{\sin x}\Bigr)_x\geq 0,\qquad
|u_x|\leq\frac{u}{\sin x}.
\end{equation}
Moreover, with $\mathcal B:=1+\lambda+u_x+u\cot x$, $\mathcal B_x\leq 0$ on
$(0,\pi)$.
\end{lemma}
\begin{proof}
Since $v$ is even about $0$ and $\pi$, $u=v_x$ is odd about both
endpoints, and in particular $u(0)=u(\pi)=0$. Differentiating~\eqref{eq:HJB} in $x$ gives
\begin{equation}\label{eq:u-eq}
-u_{xx}+\lambda u+uu_x=\zeta\sin x.
\end{equation}

We first show $u>0$ on $(0,\pi)$. Writing~\eqref{eq:u-eq} as
$Lu=\zeta\sin x$ for the linear operator
$Lw:=-w_{xx}+uw_x+\lambda w$, we have $L(0)=0<\zeta\sin x$ on $(0,\pi)$
and $u$ vanishes at both endpoints, so the strong comparison principle
gives $u>0$ on $(0,\pi)$.

We turn to $u_{xx}\leq 0$ and $(u/\sin x)_x\geq 0$, which we prove
simultaneously by tracking two quantities that record different aspects
of the shape of $u$.
The first,
\begin{equation}\label{eq:Kdef}
\mathcal K:=\frac{\zeta\sin x}{u}-\lambda-u_x ,
\end{equation}
is~\eqref{eq:u-eq} rearranged: $u_{xx}=-u\mathcal K$, so $u_{xx}\leq 0$ is
equivalent to $\mathcal K\geq 0$. The second,
\begin{equation}\label{eq:Wdef}
\mathcal W:=\sin x\,u_x-\cos x\,u=\sin^2x\,\bigl(u/\sin x\bigr)_x ,
\end{equation}
is the Wronskian of $u$ against the forcing mode $\sin x$, non-negative
precisely when $u/\sin x$ is nondecreasing.
Differentiating~\eqref{eq:Wdef} and~\eqref{eq:Kdef} and substituting
$u_{xx}=-u\mathcal K$ in both yields the coupled system
\begin{subequations}\label{eq:WK}
\begin{align}
\mathcal W_x &= \sin x\,u\,(1-\mathcal K), \label{eq:Wx}\\
\mathcal K_x &= u\mathcal K-\frac{\zeta\mathcal W}{u^2} . \label{eq:Kx}
\end{align}
\end{subequations}

Geometrically, $\mathcal K$ is the normalized curvature of $u$: the forcing
mode $\sin x$ has unit normalized curvature ($(\sin x)_{xx}=-\sin x$),
and $\mathcal K\geq 0$ records concavity of $u$. Equation~\eqref{eq:Wx} is
the standard identity for the derivative of the Wronskian, and shows that the relative slope
of $u$ against $\sin x$ moves with the curvature difference $1-\mathcal K$.
Equation~\eqref{eq:Kx} follows from $(\sin x/u)_x=-\mathcal W/u^2$, and the
system closes because the forcing is the first Dirichlet mode.

We first show $\mathcal W\geq 0$ on $(0,\pi)$. Otherwise $\mathcal W$ has a negative
interior minimum at some $x_0$. There $\mathcal W_x(x_0)=0$, so~\eqref{eq:Wx}
forces $\mathcal K(x_0)=1$. Then~\eqref{eq:Kx} gives
$\mathcal K_x(x_0)=u(x_0)-\zeta\mathcal W(x_0)/u(x_0)^2>0$.
Differentiating~\eqref{eq:Wx} once more at $x_0$,
\[
\mathcal W_{xx}(x_0)=-\sin x_0\,u(x_0)\,\mathcal K_x(x_0)<0 ,
\]
contradicting that $x_0$ is a minimum.

Next we show $\mathcal K\geq 0$ on $(0,\pi)$. At an interior zero of $\mathcal K$,
\eqref{eq:Kx} gives $\mathcal K_x=-\zeta\mathcal W/u^2\leq 0$, with equality only if
$\mathcal W=\mathcal K=0$ simultaneously. But since $\mathcal W\geq 0$ on $(0,\pi)$, an interior
point with $\mathcal W=0$ is a minimum of $\mathcal W$ and so satisfies $\mathcal W_x=0$,
while~\eqref{eq:Wx} forces $\mathcal W_x=\sin x\,u>0$ there. Hence $\mathcal K$ can cross
zero only downward.

To rule out any such crossing, we compute the limit of $\mathcal K$ at $\pi$.
Set $\xi:=\pi-x$ and expand $u=a\xi+b\xi^3+o(\xi^3)$, where the
absence of even-order terms reflects oddness of $u$ about $\pi$ and
$a:=-u_x(\pi)$. Substituting into~\eqref{eq:u-eq} gives the $\xi^1$
balance
\begin{equation}\label{eq:ab}
6b=\lambda a-a^2-\zeta .
\end{equation}
In particular $a>0$: if $a=0$ then~\eqref{eq:ab} gives $b=-\zeta/6<0$,
forcing $u\sim-\zeta\xi^3/6<0$ near $\pi$ and contradicting $u>0$ on
$(0,\pi)$. L'H\^opital's rule then gives $\sin x/u\to 1/a$ as $x\to\pi^-$, so
\[
\mathcal K(\pi^-)=\frac{\zeta}{a}+a-\lambda=\frac{a^2-\lambda a+\zeta}{a},
\]
and a direct computation using~\eqref{eq:ab} yields
\[
\mathcal W=\frac{a^2-(\lambda+1)a+\zeta}{3}\,\xi^3+o(\xi^3).
\]
The bound $\mathcal W\geq 0$ near $\pi$ forces $a^2-(\lambda+1)a+\zeta\geq 0$,
equivalently $\mathcal K(\pi^-)\geq 1>0$. By continuity $\mathcal K>0$ on a one-sided
neighborhood of $\pi$, so any $\mathcal K(x_*)<0$ at $x_*\in(0,\pi)$ would
require an upward crossing of zero on $[x_*,\pi)$, contradicting the
downward-only crossing property. Hence $\mathcal K\geq 0$ on $(0,\pi)$, i.e.\
$u_{xx}\leq 0$.

It remains to prove the chord bound and $\mathcal B_x\leq 0$.
Concavity of $u$ together with $u(0)=0$ gives $u_x\leq u/x$ wherever
$u_x>0$, and similarly $-u_x\leq u/(\pi-x)$ wherever $u_x<0$. Since
$\sin x\leq\min(x,\pi-x)$ on $(0,\pi)$,
\begin{equation}\label{eq:chord}
|u_x|\leq\frac{u}{\sin x} .
\end{equation}
For $\mathcal B_x\leq 0$, differentiate directly:
\[
\mathcal B_x=u_{xx}+\frac{u_x\sin x\cos x-u}{\sin^2x}.
\]
The first summand is $\leq 0$ by $u_{xx}\leq 0$, and the numerator of
the second satisfies $u_x\sin x\cos x\leq|u_x|\sin x\leq u$ by
$|\cos x|\leq 1$ and~\eqref{eq:chord}, so the second summand is also
$\leq 0$.
\end{proof}

The next lemma proves derivative bounds for the rescaled quotient
$Z$, used later in the change-of-variables argument for the
cubic moment (Proposition~\ref{prop:cubic}) and in the
reflection-and-correlation argument for the gradient moment
(Proposition~\ref{prop:gradient}).
Rescaling by $\sin x$ is natural because both $z$ and $\sin x$ are
smooth and odd about each endpoint of $[0,\pi]$, so $Z$ extends
smoothly to the closed interval. The proof is a chain of comparison
arguments, first for $z$, then for the flux $f\sin^2x\,Z_x$, and
finally for the flux $-(fZ)_x$.

\begin{lemma}[Bounds for $Z$]\label{lem:Z-bounds}
Let
\[
z:=v_{x\zeta},\qquad Z:=\frac{z}{\sin x}.
\]
Then $z>0$ and $Z>0$ on $(0,\pi)$, and
\[
0\leq Z_x\leq uZ\qquad\text{on }(0,\pi).
\]
The upper bound is equivalent to $(fZ)_x\leq 0$.
\end{lemma}
\begin{proof}
Differentiating~\eqref{eq:HJB} in $\zeta$ and then in $x$ produces the
equation for $z$:
\begin{equation}\label{eq:z-eq}
-z_{xx}+\lambda z+(uz)_x=\sin x,
\end{equation}
with $z(0)=z(\pi)=0$, since $z=v_{x\zeta}$ is odd at both endpoints
by Proposition~\ref{prop:smooth-zeta}.

We first show $z>0$ on $(0,\pi)$. The substitution $z=e^v y$ converts
\eqref{eq:z-eq} into
\[
-y_{xx}-uy_x+\lambda y=e^{-v}\sin x>0 ,
\]
with $y(0)=y(\pi)=0$. At an interior minimum $x_0$ of $y$,
$y_x(x_0)=0$ and $y_{xx}(x_0)\geq 0$, so the equation gives
$\lambda y(x_0)=y_{xx}(x_0)+e^{-v(x_0)}\sin x_0>0$. Hence $y>0$ on
$(0,\pi)$ and $z=e^v y>0$.

We turn to $Z_x\geq 0$. Substituting $z=Z\sin x$
into~\eqref{eq:z-eq} gives
\begin{equation}\label{eq:Z-eq}
-Z_{xx}+(u-2\cot x)Z_x+\mathcal B Z=1,
\end{equation}
which after multiplication by $f\sin^2x$ (using $f_x=-uf$) takes the flux
form
\begin{equation}\label{eq:Z-flux}
\bigl(f\sin^2x\,Z_x\bigr)_x = f\sin^2x\,(\mathcal B Z-1).
\end{equation}
Since $f\sin^2x>0$ on $(0,\pi)$, proving $Z_x\geq0$ is equivalent to
proving non-negativity of the flux $\Psi:=f\sin^2x\,Z_x$. This flux
vanishes at both endpoints. If $\Psi$ had a negative interior minimum at
$x_0$, then $\Psi_x(x_0)=0$ and
\eqref{eq:Z-flux} would give $\mathcal B(x_0)Z(x_0)=1$, hence
$\mathcal B(x_0)>0$ (since $Z>0$). Differentiating \eqref{eq:Z-flux} at
$x_0$,
\[
\Psi_{xx}(x_0)=f\sin^2x_0\bigl(\mathcal B_xZ+\mathcal B Z_x\bigr)(x_0)<0,
\]
because $\mathcal B_x\leq 0$ (Lemma~\ref{lem:u-shape}) and
$Z_x(x_0)<0$ (from $\Psi(x_0)<0$). This contradicts $x_0$ being a
minimum, so $\Psi\geq 0$, i.e.\
\begin{equation}\label{eq:Z-mono}
Z_x\geq 0.
\end{equation}

It remains to prove $Z_x\leq uZ$, which is equivalent to $\Sigma\geq 0$ for
\[
\Sigma:=uZ-Z_x=-\frac{(fZ)_x}{f}.
\]
Computing $\Sigma_x$ from~\eqref{eq:Z-eq},
\begin{equation}\label{eq:Sigma}
\Sigma_x=1-A_0Z-2\cot x\,\Sigma,\qquad A_0:=1+\lambda-u\cot x .
\end{equation}
The chord bound \eqref{eq:chord} from Lemma~\ref{lem:u-shape} gives
\[
(A_0)_x=\frac{u-u_x\sin x\cos x}{\sin^2x}\geq 0.
\]
Since $\Sigma$ vanishes at both endpoints, suppose for contradiction
that $\Sigma$ has a negative interior minimum at $x_0$. We claim
$A_0(x_0)>0$. If $x_0\in(\pi/2,\pi)$, this is immediate from
$u\cot x_0\leq 0$. If $x_0\in(0,\pi/2]$, it follows from
\eqref{eq:Sigma} at $x_0$ (where $\Sigma_x=0$, $\Sigma<0$, and
$\cot x_0\geq 0$), giving
$A_0(x_0)Z(x_0)=1-2\cot x_0\,\Sigma(x_0)\geq 1>0$. Also
$Z_x(x_0)=uZ(x_0)-\Sigma(x_0)>0$. So
$(A_0Z)_x(x_0)=(A_0)_xZ+A_0Z_x>0$. Differentiating \eqref{eq:Sigma},
\[
\Sigma_{xx}(x_0)=-(A_0Z)_x(x_0)+2\csc^2x_0\,\Sigma(x_0)<0 ,
\]
contradicting $x_0$ being a minimum. So $\Sigma\geq 0$, i.e.\
\begin{equation}\label{eq:Z-upper}
Z_x\leq uZ.
\end{equation}

Together, \eqref{eq:Z-mono} and~\eqref{eq:Z-upper} give the claimed
two-sided estimate $0\leq Z_x\leq uZ$. In addition, since $f_x=-uf$,
the upper bound gives
\[
(fZ)_x=f(Z_x-uZ)\leq0.
\]
\end{proof}

\section{Proof of Theorem~\ref{thm:main}}\label{sec:proof}

The proof of part~(i) (in its dimensionless form $A''(\zeta)<0$)
has three steps. We first establish an exact identity decomposing
$A''(\zeta)$ as the sum of a cubic moment and a gradient moment
(Lemma~\ref{lem:second-deriv}). We then show that the cubic moment is
nonnegative (Proposition~\ref{prop:cubic}) and the gradient moment is
strictly positive (Proposition~\ref{prop:gradient}). The two parts of the
theorem are then assembled at the end of the section.

\begin{lemma}[Second-derivative identity]\label{lem:second-deriv}
Let
\[
w:=\frac{A}{\lambda}+v_\zeta,\qquad z:=v_{x\zeta}.
\]
Then
\[
\int_\T wf=0,\qquad
A''(\zeta)=-\lambda\int_\T w^3f-3\int_\T wz^2f .
\]
Equivalently, with $X\sim\mu^\zeta$, $W:=w(X)$, and $Y:=z(X)$,
\[
\mathbb E[W]=0,\qquad
A''(\zeta)=-\lambda\,\mathbb E[W^3]-3\,\mathbb E[WY^2].
\]
\end{lemma}
\begin{proof}
Differentiating~\eqref{eq:HJB} in $\zeta$, multiplying by $f$, and
using $f_x=-uf$ gives the divergence form
\begin{equation}\label{eq:v-zeta-eq}
-(fv_{x\zeta})_x+\lambda fv_\zeta=-f\cos x.
\end{equation}
Integrating over $\T$ yields $\lambda\int_\T v_\zeta f=-A$, hence
$\int_\T wf=0$. Combined with $f=e^{-v^\zeta}/\mathcal Z^\zeta$, this gives
\begin{equation}\label{eq:f-zeta}
f_\zeta=-wf,
\end{equation}
since $f_\zeta=-v_\zeta f-(\mathcal Z^\zeta_\zeta/\mathcal Z^\zeta)f$ and
$\mathcal Z^\zeta_\zeta/\mathcal Z^\zeta=-\int_\T v_\zeta f=A/\lambda$.

Rewriting~\eqref{eq:v-zeta-eq} in terms of $w=A/\lambda+v_\zeta$,
\begin{equation}\label{eq:w-energy-eq}
-(fw_x)_x+\lambda fw=(A-\cos x)f .
\end{equation}
Differentiating $A=\int_\T\cos x\,f$ and using~\eqref{eq:f-zeta}
together with $\int_\T wf=0$,
\[
A'=-\int_\T \cos x\,wf=\int_\T (A-\cos x)wf .
\]
Testing~\eqref{eq:w-energy-eq} against $w$ then yields the energy
identity
\begin{equation}\label{eq:A-prime-energy}
A'=\int_\T z^2f+\lambda\int_\T w^2f .
\end{equation}

Differentiating~\eqref{eq:HJB} twice in $\zeta$ and multiplying by
$f$, the same divergence-form computation that produced
\eqref{eq:v-zeta-eq} gives, with $b:=w_\zeta$ (so that $b_x=z_\zeta$),
\begin{equation}\label{eq:b-eq}
-(fb_x)_x+\lambda fb=A'f-z^2f .
\end{equation}
Testing~\eqref{eq:b-eq} against $w$ and using $\int_\T wf=0$,
\begin{equation}\label{eq:b-tested}
\int_\T zz_\zeta f+\lambda\int_\T ww_\zeta f
=-\int_\T wz^2f .
\end{equation}
Finally, differentiating~\eqref{eq:A-prime-energy} in $\zeta$ and
applying~\eqref{eq:f-zeta},
\[
A''=2\int_\T zz_\zeta f+2\lambda\int_\T ww_\zeta f
-\int_\T wz^2f-\lambda\int_\T w^3f,
\]
and substituting~\eqref{eq:b-tested} gives the claimed formula.
\end{proof}

The cubic moment $\int_\T w^3 f$ in the identity above is handled in
Proposition~\ref{prop:cubic} via a change of variables
$C(x)=\cos\theta(x)$ under which $w$ becomes affine in $\cos\theta$
and $f$ pushes forward to an explicit density $\eta$ on $[0,\pi]$. The
two ingredients are a shape statement on $v_\zeta$, which will force
$\eta$ to be nonincreasing (Lemma~\ref{lem:geometric-mean} below), and
an abstract inequality for centered cubic moments of $\cos\theta$
against nonincreasing densities (Lemma~\ref{lem:monotone-cosine-skewness}).

\begin{lemma}[Geometric-mean monotonicity]\label{lem:geometric-mean}
Let
\[
a:=v_\zeta,\qquad
\Delta_0(x):=a(x)-a(0),\qquad \Delta_\pi(x):=a(\pi)-a(x).
\]
Then
\[
H(x):=e^{-v(x)}\frac{\sqrt{\Delta_0(x)\Delta_\pi(x)}}{\sin x}
\]
is nonincreasing on $(0,\pi)$.
\end{lemma}
\begin{proof}
With $z, Z$ as in Lemma~\ref{lem:Z-bounds} and $z=a_x$,
$z,Z>0$ on $(0,\pi)$ and $(e^{-v}Z)_x\leq 0$.

To identify pointwise bounds on $\Delta_0,\Delta_\pi$ that force $H_x\leq 0$, use
$\sin^2 x=(1-\cos x)(1+\cos x)$ to factor
\[
H^2=\Bigl(e^{-v}\,\frac{\Delta_0}{1-\cos x}\Bigr)\Bigl(e^{-v}\,\frac{\Delta_\pi}{1+\cos x}\Bigr),
\]
so that $2\log H$ is the sum of the two logarithms. Using $(\Delta_0)_x=z$
and $(\Delta_\pi)_x=-z$,
\[
\partial_x\log\Bigl(e^{-v}\frac{\Delta_0}{1-\cos x}\Bigr)=-u+\frac z{\Delta_0}-\frac{\sin x}{1-\cos x},\qquad
\partial_x\log\Bigl(e^{-v}\frac{\Delta_\pi}{1+\cos x}\Bigr)=-u-\frac z{\Delta_\pi}+\frac{\sin x}{1+\cos x}.
\]
The trigonometric remainders are half-angle quantities. With
\begin{equation}\label{eq:rt-def}
r:=\cot\tfrac x2=\frac{\sin x}{1-\cos x}=\frac{1+\cos x}{\sin x},\qquad
t:=\tan\tfrac x2=\frac{\sin x}{1+\cos x}=\frac{1-\cos x}{\sin x},
\end{equation}
which satisfy $rt=1$, the log-derivatives simplify to
\[
\partial_x\log\Bigl(e^{-v}\frac{\Delta_0}{1-\cos x}\Bigr)=-u+\frac{z}{\Delta_0}-r,\qquad
\partial_x\log\Bigl(e^{-v}\frac{\Delta_\pi}{1+\cos x}\Bigr)=-u-\frac{z}{\Delta_\pi}+t.
\]
The first is $\leq 0$ iff $z/\Delta_0\leq u+r$, i.e.\ $\Delta_0\geq z/(u+r)$.
The second is automatically $\leq 0$ where $t\leq u$ (both $-u+t$ and
$-z/\Delta_\pi$ are nonpositive), and where $t>u$ it reduces to
$z/\Delta_\pi\geq t-u$, i.e.\ $\Delta_\pi\leq z/(t-u)$. So it suffices to establish
\begin{equation}\label{eq:a-pm-bounds}
\Delta_0(x)\geq\frac{z(x)}{u(x)+r(x)},\qquad
\Delta_\pi(x)\leq\frac{z(x)}{t(x)-u(x)}\ \text{on $\{t>u\}$.}
\end{equation}

We translate these into integrals via the monotonicity of $e^{-v}Z$.
Since $(e^{-v}Z)_x\leq 0$,
\[
Z(y)\geq e^{v(y)-v(x)}Z(x)\text{ for }y\leq x,\qquad
Z(y)\leq e^{v(y)-v(x)}Z(x)\text{ for }y\geq x,
\]
and, since $a_x=z=Z\sin x$, multiplying by $\sin y$ and integrating gives
\begin{equation}\label{eq:a-ZI}
\Delta_0(x)\geq Z(x)\,I_0(x),\qquad \Delta_\pi(x)\leq Z(x)\,I_\pi(x),
\end{equation}
where
\[
I_0(x):=e^{-v(x)}\int_0^x e^{v(y)}\sin y\,dy,\qquad
I_\pi(x):=e^{-v(x)}\int_x^\pi e^{v(y)}\sin y\,dy
\]
solve, by differentiating under the integral and using $v_x=u$, the
first-order ODEs
\begin{equation}\label{eq:I-odes}
I_0'+uI_0=\sin x,\quad I_0(0)=0;\qquad
I_\pi'+uI_\pi=-\sin x,\quad I_\pi(\pi)=0.
\end{equation}
By~\eqref{eq:a-ZI} and $Z=z/\sin x$, the bounds~\eqref{eq:a-pm-bounds} will follow from
\begin{equation}\label{eq:I-bounds}
I_0(x)\geq\frac{\sin x}{u(x)+r(x)},\qquad
I_\pi(x)\leq\frac{\sin x}{t(x)-u(x)}\ \text{on $\{t>u\}$,}
\end{equation}
which we now prove. Both proofs use the derivative identities
\begin{equation}\label{eq:half-angle-diff}
\sin x\cdot r'=-r,\qquad \sin x\cdot t'=t,
\end{equation}
obtained by differentiating either form in~\eqref{eq:rt-def}.

For the lower bound on $I_0$, set
\[
D_0:=\frac{\sin x}{u+r}.
\]
Differentiating
$D_0(u+r)=\sin x$ and using $r'=-r/\sin x$
from~\eqref{eq:half-angle-diff} and $r\sin x=1+\cos x$
from~\eqref{eq:rt-def} to simplify the resulting algebra yields
\begin{equation}\label{eq:D0-ode}
D_0'+uD_0=\sin x-\frac{u+\sin x\,u_x}{(u+r)^2}.
\end{equation}
The chord bound $|u_x|\leq u/\sin x$ from Lemma~\ref{lem:u-shape} gives
$\sin x\,u_x\geq -u$, so the subtracted term in~\eqref{eq:D0-ode} is
nonnegative. Subtracting~\eqref{eq:D0-ode} from~\eqref{eq:I-odes} and
multiplying by $e^v$ gives
\[
\bigl(e^v(I_0-D_0)\bigr)_x=e^v\bigl[(I_0-D_0)'+u(I_0-D_0)\bigr]\geq 0.
\]
As $x\to 0^+$, $u\to 0$ and $r\sim 2/x$, so $D_0\sim x^2/2$. Likewise
$I_0\sim\int_0^x y\,dy=x^2/2$. Hence $e^v(I_0-D_0)$ vanishes at $0$ and
is nondecreasing, giving $I_0\geq D_0$.

For the upper bound on $I_\pi$, first observe that $\{t>u\}$ is an
interval $(x_*,\pi)$ for some $x_*\in[0,\pi)$. Indeed, from $t'=t/\sin x$
in~\eqref{eq:half-angle-diff} and the chord bound~\eqref{eq:chord},
\[
\Bigl(\frac u t\Bigr)_x=\frac{u_x-u/\sin x}{t}\leq 0,
\]
so $u/t$ is nonincreasing on $(0,\pi)$. Since $u(\pi)=0$ and
$t(\pi)=+\infty$, $u/t\to 0$ at $\pi$, so $\{u/t<1\}=\{t>u\}$ is a
right-neighborhood of $\pi$, and it has the form
$\{t>u\}=(x_*,\pi)$. On this interval, set
\[
D_\pi:=\frac{\sin x}{t-u}.
\]
Differentiating $D_\pi(t-u)=\sin x$ and using $\sin x\cdot t'=t$
together with $t\sin x=1-\cos x$ to simplify gives
\[
D_\pi'+uD_\pi=-\sin x+\frac{\sin x\,u_x-u}{(t-u)^2}.
\]
The chord bound~\eqref{eq:chord} gives $\sin x\,u_x\leq u$, so the
last term on the right is nonpositive. Subtracting this from~\eqref{eq:I-odes} and multiplying
by $e^v$ gives
\[
\bigl(e^v(D_\pi-I_\pi)\bigr)_x=e^v\bigl[(D_\pi-I_\pi)'+u(D_\pi-I_\pi)\bigr]\leq 0.
\]
At $x=\pi^-$ both $I_\pi$ and $D_\pi$ tend to $0$ ($I_\pi$ from its
definition, $D_\pi$ because $\sin x\to 0$ while $t-u\to\infty$).
Integrating backward from $\pi$ yields $D_\pi\geq I_\pi$.

This establishes~\eqref{eq:I-bounds}, hence~\eqref{eq:a-pm-bounds},
hence $2(\log H)_x\leq 0$, so $H$ is nonincreasing.
\end{proof}

The next lemma is the abstract counterpart to the geometric-mean
monotonicity just established: for any nonincreasing probability density
on $[0,\pi]$, the centered cubic moment of $\cos\theta$ has a definite
sign. It is purely measure-theoretic and makes no reference to the mean
field game.

\begin{lemma}[Monotone cosine skewness]\label{lem:monotone-cosine-skewness}
Let $\eta\geq0$ be a nonincreasing probability density on $[0,\pi]$ and set
\[
\bar c:=\int_0^\pi \cos\theta\,\eta(\theta)\,d\theta .
\]
Then
\[
\int_0^\pi(\bar c-\cos\theta)^3\eta(\theta)\,d\theta\geq 0.
\]
\end{lemma}
\begin{proof}
On $[0,\pi/2]$ monotonicity gives $\eta(\theta)\geq\eta(\pi-\theta)$. Combined with $\cos(\pi-\theta)=-\cos\theta$ this yields
\[
\bar c=\int_0^{\pi/2}\cos\theta\,\bigl(\eta(\theta)-\eta(\pi-\theta)\bigr)\,d\theta\geq 0.
\]
By definition of $\bar c$,
\begin{equation}\label{eq:zero-mean-cos}
\int_0^\pi(\cos\theta-\bar c)\eta(\theta)\,d\theta=0.
\end{equation}
We prove the equivalent inequality
\[
\int_0^\pi(\cos\theta-\bar c)^3\eta(\theta)\,d\theta\leq 0.
\]

Let $\theta_0:=\arccos\bar c$. On $[0,\theta_0]$, $\cos\theta$ decreases
from $1$ to $\bar c$, so setting $\tau:=\cos\theta-\bar c$ gives
$\tau\in[0,1-\bar c]$ with
$d\tau=-\sin\theta\,d\theta=-\sqrt{1-(\bar c+\tau)^2}\,d\theta$, hence
pushforward density
\[
\alpha_+(\tau):=\frac{\eta(\arccos(\bar c+\tau))}{\sqrt{1-(\bar c+\tau)^2}},\qquad\tau\in[0,1-\bar c].
\]
On $[\theta_0,\pi]$, $\cos\theta$ decreases from $\bar c$ to $-1$, so
setting $\tau:=\bar c-\cos\theta$ gives $\tau\in[0,1+\bar c]$ and the
analogous computation yields
\[
\alpha_-(\tau):=\frac{\eta(\arccos(\bar c-\tau))}{\sqrt{1-(\bar c-\tau)^2}},\qquad\tau\in[0,1+\bar c].
\]
The supports differ ($1+\bar c\geq 1-\bar c$ since $\bar c\geq 0$), but
on the common subinterval $[0,1-\bar c]$ we have $\alpha_+\geq\alpha_-$:
since $\bar c+\tau\geq\bar c-\tau$ and $\arccos$ is decreasing,
monotonicity of $\eta$ gives the numerator inequality
$\eta(\arccos(\bar c+\tau))\geq\eta(\arccos(\bar c-\tau))$, and
$(\bar c+\tau)^2\geq(\bar c-\tau)^2$ (using $\bar c\geq 0$) gives the
denominator one.

In these variables, \eqref{eq:zero-mean-cos} reads
\begin{equation}\label{eq:zm-tau}
\int_0^{1-\bar c}\tau\,(\alpha_+-\alpha_-)\,d\tau
=\int_{1-\bar c}^{1+\bar c}\tau\,\alpha_-(\tau)\,d\tau ,
\end{equation}
and the centered cubic moment becomes
\[
\int_0^\pi(\cos\theta-\bar c)^3\eta\,d\theta
=\int_0^{1-\bar c}\tau^3(\alpha_+-\alpha_-)\,d\tau
-\int_{1-\bar c}^{1+\bar c}\tau^3\,\alpha_-(\tau)\,d\tau .
\]
Since $\tau^2\leq(1-\bar c)^2$ on $[0,1-\bar c]$ and
$\tau^2\geq(1-\bar c)^2$ on $[1-\bar c,1+\bar c]$,
\[
\int_0^{1-\bar c}\tau^3(\alpha_+-\alpha_-)\,d\tau
\leq (1-\bar c)^2\int_0^{1-\bar c}\tau(\alpha_+-\alpha_-)\,d\tau
\stackrel{\eqref{eq:zm-tau}}{=}(1-\bar c)^2\int_{1-\bar c}^{1+\bar c}\tau\,\alpha_-\,d\tau
\leq \int_{1-\bar c}^{1+\bar c}\tau^3\,\alpha_-\,d\tau .
\]
Hence $\int_0^\pi(\cos\theta-\bar c)^3\eta\,d\theta\leq 0$, which is
the stated inequality.
\end{proof}

With both ingredients in hand we prove that the cubic moment is nonnegative.
The geometric-mean monotonicity (Lemma~\ref{lem:geometric-mean}) is
used to show that the pushforward density $\eta$ is nonincreasing, and
the cosine-skewness inequality (Lemma~\ref{lem:monotone-cosine-skewness})
then yields the sign.

\begin{proposition}[Sign of the cubic moment]\label{prop:cubic}
For every $\zeta>0$,
\[
\int_\T w^3 f\geq 0.
\]
\end{proposition}
\begin{proof}
With $a:=v_\zeta$ and $z, Z$ as in Lemma~\ref{lem:Z-bounds} (so $z=a_x$), set
\[
\Delta:=a(\pi)-a(0),\qquad
\Delta_0(x):=a(x)-a(0),\qquad \Delta_\pi(x):=a(\pi)-a(x),
\]
\[
C:=\frac{\Delta_\pi-\Delta_0}{\Delta},\qquad S:=1-C^2,
\]
so that $w=A/\lambda+a$ in the notation of Lemma~\ref{lem:second-deriv}.
By Lemma~\ref{lem:Z-bounds}, $z>0$ on $(0,\pi)$, so $\Delta>0$ and
\begin{equation}\label{eq:Cx-formula}
C_x=-\frac{2z}{\Delta}=-\frac{2}{\Delta}\sin x\,Z<0.
\end{equation}
Hence $C\colon[0,\pi]\to[-1,1]$ is a strictly
decreasing bijection from $C(0)=1$ to $C(\pi)=-1$, and we define
$\theta\colon[0,\pi]\to[0,\pi]$ by
\[
C(x)=\cos\theta(x).
\]
Let
\[
\eta(\theta):=\frac{2f(x(\theta))}{\theta'(x(\theta))}
\]
be the pushforward density of $2f(x)\,dx$ under this map. Its total mass is
\[
\int_0^\pi \eta(\theta)\,d\theta
=\int_0^\pi \frac{2f(x(\theta))}{\theta'(x(\theta))}\,d\theta
=2\int_0^\pi f(x)\,dx
=\int_\T f(x)\,dx=1,
\]
where the penultimate equality uses the evenness of $f$. Thus $\eta$ is
a probability density on $[0,\pi]$.

We first express the cubic moment as a centered cubic moment of
$\cos\theta$ under $\eta$. Since
$\Delta_\pi-\Delta_0=a(\pi)+a(0)-2a$ and $C=(\Delta_\pi-\Delta_0)/\Delta$, $a$ is affine in
$C$:
\[
a(x)=\frac{a(\pi)+a(0)}{2}-\frac{\Delta}{2}\,C(x) .
\]
The zero-mean identity $\int_\T wf=0$ from
Lemma~\ref{lem:second-deriv} fixes the additive constant and gives
\begin{equation}\label{eq:affine-C}
w(x)=\frac{\Delta}{2}\bigl(\bar c-C(x)\bigr),\qquad
\bar c:=\int_0^\pi\cos\theta\,\eta(\theta)\,d\theta .
\end{equation}
Cubing \eqref{eq:affine-C} and pushing forward to $\theta$ gives
\[
\int_\T w^3 f\,dx
=\frac{\Delta^3}{8}\int_0^\pi(\bar c-\cos\theta)^3\eta(\theta)\,d\theta .
\]
By Lemma~\ref{lem:monotone-cosine-skewness}, the right-hand side is
nonnegative provided $\eta$ is nonincreasing. It therefore remains to
prove that $\eta$ is nonincreasing.

From $C=\cos\theta$ and $C_x<0$ we get
$\theta'=-C_x/\sin\theta=-C_x/\sqrt S$. Differentiating
$\log\eta=\log(2f)-\log\theta'$ in $\theta$ and using $(\log f)_x=-u$,
\[
\frac{d}{d\theta}\log\eta=-\frac{1}{\theta'}\left(u+\frac{\theta''}{\theta'}\right),
\]
which, after substituting $\theta'=-C_x/\sqrt S$, gives that $\eta$ is
nonincreasing iff
\begin{equation}\label{eq:eta-monotone-condition}
S\left(u+\frac{C_{xx}}{C_x}\right)+CC_x\geq0.
\end{equation}
Logarithmic differentiation of \eqref{eq:Cx-formula} gives
$C_{xx}/C_x=\cot x+Z_x/Z$, so the left side of
\eqref{eq:eta-monotone-condition} splits as
\begin{equation}\label{eq:eta-split}
S\left(u+\frac{C_{xx}}{C_x}\right)+CC_x
=\bigl[S(u+\cot x)+CC_x\bigr]+S\frac{Z_x}{Z}.
\end{equation}
The second term is nonnegative since $S=1-C^2\geq 0$ and $Z_x\geq 0$ by
Lemma~\ref{lem:Z-bounds}. For the first term, the geometric-mean
monotonicity (Lemma~\ref{lem:geometric-mean}) gives
\[
\frac{H_x}{H}
=-\left(u+\cot x-\frac{z(\Delta_\pi-\Delta_0)}{2\Delta_0\Delta_\pi}\right)\leq 0,
\]
and using $S=4\Delta_0\Delta_\pi/\Delta^2$ (since $\Delta_\pi+\Delta_0=\Delta$) together with
\eqref{eq:Cx-formula}, which gives $CC_x=-2z(\Delta_\pi-\Delta_0)/\Delta^2$, one
identifies
\[
S(u+\cot x)+CC_x=-S\frac{H_x}{H}\geq 0.
\]
Hence \eqref{eq:eta-monotone-condition} holds, and $\eta$ is
nonincreasing.
\end{proof}

The estimate for the gradient moment relies on the following central-positivity statement.

\begin{lemma}[Central positivity]\label{lem:central-positivity}
Assume $\zeta>0$. Let $\Omega\in C^1([0,\pi])$ be non-constant and
satisfy
\[
\Omega(\pi-x)=\Omega(x),\qquad \Omega'(x)\geq 0\quad\text{for }0<x<\pi/2.
\]
Then, with $w$ as in Lemma~\ref{lem:second-deriv},
\begin{equation}\label{eq:omega-positive}
\int_0^\pi w(x)\,\Omega(x)\,f(x)\,dx> 0 .
\end{equation}
\end{lemma}
\begin{proof}
Set $\bar x:=\pi-x$. With $z$ as in Lemma~\ref{lem:Z-bounds},
$w_x=z>0$, so $w$ is strictly increasing on $(0,\pi)$. Since $wf$ is
even about $0$ (Proposition~\ref{prop:smooth-zeta}),
Lemma~\ref{lem:second-deriv} gives $\int_0^\pi wf=0$.

The strategy is to recover $\Omega$ in~\eqref{eq:omega-positive} from a
family of central-window estimates via the layer-cake-style decomposition
\begin{equation}\label{eq:omega-decomp}
\Omega(x)=\Omega(0)+\int_0^{\pi/2}\Omega'(\alpha)\,
\mathbf 1_{\{\alpha\leq x\leq\pi-\alpha\}}\,d\alpha,
\end{equation}
which holds for every $C^1$ function on $[0,\pi]$ symmetric about
$\pi/2$ by the fundamental theorem of calculus.
Substituting~\eqref{eq:omega-decomp} into~\eqref{eq:omega-positive}
and using $\int_0^\pi wf=0$ to discard the $\Omega(0)$ term yields
\[
\int_0^\pi w(x)\,\Omega(x)\,f(x)\,dx
=\int_0^{\pi/2}\Omega'(\alpha)\!\int_\alpha^{\pi-\alpha}w(x)f(x)\,dx\,d\alpha.
\]
Since $\Omega'\geq 0$ on $(0,\pi/2)$ with strict positivity on a set of
positive measure ($\Omega$ is $C^1$ and non-constant),
\eqref{eq:omega-positive} reduces to
\begin{equation}\label{eq:central-positive}
\int_\alpha^{\pi-\alpha}w(x)f(x)\,dx>0,\qquad 0<\alpha<\pi/2.
\end{equation}

Set
\[
\rho(x):=\frac{f(\bar x)}{f(x)},\qquad 0\leq x\leq \pi/2.
\]
Since $f$ is strictly decreasing on $(0,\pi)$ (because $f_x=-uf$ with
$u>0$), $\rho$ is strictly increasing on $[0,\pi/2]$, with
$\rho(\pi/2)=1$ and $\rho'>0$ on $(0,\pi/2)$. The change of
variables $y=\bar x$ gives the reflection identity
\begin{equation}\label{eq:reflect-identity}
\int_0^\alpha \rho(x)f(x)\,g(\bar x)\,dx=\int_{\pi-\alpha}^\pi g(y)f(y)\,dy,
\qquad 0\leq\alpha\leq\pi/2,
\end{equation}
valid for any integrable $g$. By Lemma~\ref{lem:Z-bounds}, $fZ$ is
nonincreasing, so
\[
f(x)Z(x)\geq f(\bar x)Z(\bar x),\qquad 0\leq x\leq\pi/2.
\]
Multiplying by $\sin x=\sin\bar x$ and using $z=Z\sin x$,
\begin{equation}\label{eq:fz-reflect}
f(x)z(x)\geq f(\bar x)z(\bar x),\quad\text{i.e.}\quad
z(x)\geq\rho(x)z(\bar x),\qquad 0\leq x\leq \pi/2.
\end{equation}

We introduce the folded profile
\[
\tilde w(x):=w(x)+\rho(x)\,w(\bar x),\qquad 0\leq x\leq \pi/2.
\]
Applying~\eqref{eq:reflect-identity} to $g=w$ gives, for
$0\leq\alpha\leq\pi/2$,
\begin{equation}\label{eq:wt-fold}
\int_0^\alpha \tilde w(x)f(x)\,dx
=\int_0^\alpha wf+\int_{\pi-\alpha}^\pi wf .
\end{equation}
At $\alpha=\pi/2$, the right-hand side equals $\int_0^\pi wf=0$, so
\begin{equation}\label{eq:wt-zero-mean}
\int_0^{\pi/2} \tilde w(x)f(x)\,dx=0 .
\end{equation}

We first show $w(\pi/2)>0$. Fix $x\in[0,\pi/2)$.
Combining~\eqref{eq:fz-reflect} with the monotonicity of $\rho$, for
$s\in[x,\pi/2]$,
\[
z(s)\geq \rho(s)z(\bar s)\geq \rho(x)z(\bar s),
\]
the second inequality being strict on $(x,\pi/2)$. Since $w_x=z$ and
$(d/ds)w(\bar s)=-z(\bar s)$, integrating in $s$ over $[x,\pi/2]$
yields $w(\pi/2)-w(x)>\rho(x)\bigl(w(\bar x)-w(\pi/2)\bigr)$,
equivalently
\begin{equation}\label{eq:wt-below-mid}
\tilde w(x)<(1+\rho(x))\,w(\pi/2),\qquad 0\leq x<\pi/2.
\end{equation}
Integrating~\eqref{eq:wt-below-mid} against $f(x)$ over $[0,\pi/2]$,
the left-hand side is $0$ by~\eqref{eq:wt-zero-mean}, while the
right-hand side equals $w(\pi/2)\int_0^\pi f$. Since $\int_0^\pi f>0$,
$w(\pi/2)>0$.

Differentiating $\tilde w$,
\begin{equation}\label{eq:wt-mono}
\tilde w'(x)=\bigl[z(x)-\rho(x)z(\bar x)\bigr]+\rho'(x)\,w(\bar x)>0
\qquad\text{on }(0,\pi/2),
\end{equation}
since the bracket is nonnegative by~\eqref{eq:fz-reflect}, $\rho'>0$,
and $w(\bar x)\geq w(\pi/2)>0$ (as $w$ is strictly increasing on
$(0,\pi)$ and $\bar x\geq \pi/2$).

By~\eqref{eq:wt-zero-mean} and the strict
monotonicity~\eqref{eq:wt-mono},
\[
\int_0^\alpha \tilde w(x)f(x)\,dx<0,\qquad 0<\alpha<\pi/2.
\]
By~\eqref{eq:wt-fold} together with $\int_0^\pi wf=0$,
\[
\int_\alpha^{\pi-\alpha}wf
=-\int_0^\alpha wf-\int_{\pi-\alpha}^\pi wf
=-\int_0^\alpha \tilde w f>0.
\]
\end{proof}


\begin{proposition}[Sign of the gradient moment]\label{prop:gradient}
For every $\zeta>0$,
\[
\int_\T w\,z^2 f> 0.
\]
\end{proposition}
\begin{proof}
Let $Z$ be as in Lemma~\ref{lem:Z-bounds}. Since $w$, $z^2$, and $f$
are even with respect to $0$ (Proposition~\ref{prop:smooth-zeta}), it suffices
to prove the sign on $[0,\pi]$.

Let $d\nu=\sin^2x\,f(x)\,dx$ on $[0,\pi]$. Since $w_x=z$, both $w$
and $Z^2$ are nondecreasing by Lemma~\ref{lem:Z-bounds}. For any pair
of functions $p,q$ on $[0,\pi]$ one has the symmetrization identity
\[
\nu([0,\pi])\int pq\,d\nu-\int p\,d\nu\int q\,d\nu
=\frac12\iint_{[0,\pi]^2}(p(x)-p(y))(q(x)-q(y))\,d\nu(x)d\nu(y),
\]
whose right-hand side is non-negative whenever $p$ and $q$ are both
nondecreasing, since then $(p(x)-p(y))(q(x)-q(y))\geq 0$ for all $x,y$. Applying it with $p=w$ and $q=Z^2$
gives
\[
\int_0^\pi wZ^2\,d\nu
\geq
\frac{1}{\nu([0,\pi])}\int_0^\pi w\,d\nu\int_0^\pi Z^2\,d\nu .
\]
Here $\int w\,d\nu>0$ by Lemma~\ref{lem:central-positivity} applied to
$\Omega(x)=\sin^2 x$, while $\int Z^2\,d\nu>0$ because $Z>0$ on
$(0,\pi)$. Hence
\[
\int_0^\pi w(x)z(x)^2f(x)\,dx
=\int_0^\pi w(x)Z(x)^2\sin^2x\,f(x)\,dx>0.
\]
By evenness, $\int_\T wz^2 f=2\int_0^\pi wz^2 f>0$.
\end{proof}

\begin{proof}[Proof of Theorem~\ref{thm:main}]
By Lemma~\ref{lem:second-deriv},
\[
A''(\zeta)
=-\lambda\int_\T w^3 f
-3\int_\T w\,z^2 f.
\]
The first integral is nonnegative by Proposition~\ref{prop:cubic} and
the second is strictly positive by Proposition~\ref{prop:gradient}.
Since $\lambda>0$, $A''(\zeta)<0$. Hence $A$, and so the original map
$F_\kappa(\gamma)=\kappa A(4\gamma/\sigma^4)$, is strictly concave.
This is part~(i).

For part~(ii), recall the properties~\eqref{eq:F-properties} and
that, by \cite[Prop.~6.4]{CCS}, a probability measure is a
non-uniform stationary equilibrium of the Kuramoto mean field game if
and only if, up to a rotation of $\T$, it equals $\mu^{\zeta_*}$ for
some positive fixed point $\gamma_*$ of $F_\kappa$, where
$\zeta_*:=4\gamma_*/\sigma^4$ is the corresponding dimensionless
parameter. Setting
$g_\kappa(\gamma):=F_\kappa(\gamma)-\gamma$, it therefore suffices to
show that $g_\kappa$ has no positive zero for $\kappa\leq\kappa_c$ and
a unique positive zero $\gamma_*(\kappa)$ for $\kappa>\kappa_c$, that
$\kappa\mapsto\gamma_*(\kappa)$ is continuous on $(\kappa_c,\infty)$,
and that $\gamma_*(\kappa)\to 0$ as $\kappa\downarrow\kappa_c$.

From~\eqref{eq:F-properties}, $g_\kappa(0)=0$,
$g_\kappa'(0)=\kappa/\kappa_c-1$, and
$g_\kappa(\gamma)\leq\kappa-\gamma\to-\infty$ as $\gamma\to\infty$.
Part~(i) gives that $g_\kappa$ is strictly concave on $[0,\infty)$. By
strict concavity the tangent at $0$ strictly dominates:
\begin{equation}\label{eq:tangent-bound}
g_\kappa(\gamma)<g_\kappa'(0)\,\gamma\qquad\text{for }\gamma>0.
\end{equation} Hence if
$\kappa\leq\kappa_c$, then $g_\kappa'(0)\leq 0$ and $g_\kappa<0$ on
$(0,\infty)$, so there is no positive zero. If $\kappa>\kappa_c$, then
$g_\kappa'(0)>0$ makes $g_\kappa$ positive on a right neighborhood of
$0$, while $g_\kappa\to-\infty$, so by continuity $g_\kappa$ has a
positive zero. Strict concavity allows $g_\kappa$ to take any value at
most twice, and $g_\kappa(0)=0$ already accounts for one zero, so the
positive zero is unique. Denote it by $\gamma_*(\kappa)$, with
corresponding dimensionless parameter $\zeta_*(\kappa)$. Then $g_\kappa>0$ on $(0,\gamma_*(\kappa))$ and $g_\kappa<0$
on $(\gamma_*(\kappa),\infty)$, and by \cite[Prop.~6.4]{CCS},
$\mu_\kappa=\mu^{\zeta_*(\kappa)}$.

At $\gamma_*(\kappa)$, $g_\kappa$ crosses from positive to negative,
so $g_\kappa'(\gamma_*(\kappa))<0$ by strict concavity. Since
$(\kappa,\gamma)\mapsto g_\kappa(\gamma)$ is jointly smooth, the
implicit function theorem yields that $\kappa\mapsto\zeta_*(\kappa)$
is $C^\infty$ on $(\kappa_c,\infty)$. By
Proposition~\ref{prop:smooth-zeta}, $\zeta\mapsto f^\zeta$ is smooth
as a map into $C^\infty(\T)$, so the composition
$\kappa\mapsto f^{\zeta_*(\kappa)}$ is $C^\infty$ into $C^\infty(\T)$.

It remains to show $f^{\zeta_*(\kappa)}\to 1/(2\pi)$ in $C^\infty(\T)$
as $\kappa\downarrow\kappa_c$. Since $\zeta\mapsto f^\zeta$ is
continuous at $\zeta=0$ into $C^\infty(\T)$ and $f^0\equiv 1/(2\pi)$,
it suffices to show $\zeta_*(\kappa)\to 0$, equivalently
$\gamma_*(\kappa)\to 0$. Evaluating~\eqref{eq:tangent-bound} at
$\kappa=\kappa_c$ gives
$g_{\kappa_c}<0$ on $(0,\infty)$, so by joint continuity of $g_\kappa$
in $(\kappa,\gamma)$, for any $\epsilon>0$ there is $\delta>0$ with
$g_\kappa(\epsilon)<0$ for $\kappa\in(\kappa_c,\kappa_c+\delta)$. Since
$g_\kappa>0$ on $(0,\gamma_*(\kappa))$, this forces
$\gamma_*(\kappa)<\epsilon$, and $\gamma_*(\kappa)\to 0$ follows.
\end{proof}

\begin{thebibliography}{ABPVS05}

\bibitem[ABPVS05]{ABPVS}
J.~A.~Acebr\'on, L.~L.~Bonilla, C.~J.~P\'erez Vicente, F.~Ritort, and
  R.~Spigler.
\newblock The {K}uramoto model: a simple paradigm for synchronization
  phenomena.
\newblock \emph{Reviews of Modern Physics} 77 (2005), no.\ 1, 137--185.

\bibitem[BGP14]{BGP}
L.~Bertini, G.~Giacomin, and C.~Poquet.
\newblock Synchronization and random long time dynamics for mean-field
  plane rotators.
\newblock \emph{Probability Theory and Related Fields} 160 (2014),
  no.\ 3--4, 593--653.

\bibitem[CD18]{CarmonaDelarue}
R.~Carmona and F.~Delarue.
\newblock \emph{Probabilistic Theory of Mean Field Games with
  Applications, I--II}.
\newblock Probability Theory and Stochastic Modelling, vols.\ 83--84.
  Springer, 2018.

\bibitem[CCS23]{CCS}
R.~Carmona, Q.~Cormier, and H.~M.~Soner.
\newblock Synchronization in a {K}uramoto mean field game.
\newblock \emph{Communications in Partial Differential Equations}
  48 (2023), no.\ 9, 1214--1244. arXiv:2210.12912.

\bibitem[CCS25]{CCS2}
R.~Carmona, Q.~Cormier, and H.~M.~Soner.
\newblock {K}uramoto mean field game with intrinsic frequencies.
\newblock Preprint, 2025. arXiv:2509.18000.

\bibitem[CC24]{CesaroniCirant}
A.~Cesaroni and M.~Cirant.
\newblock Stationary equilibria and their stability in a {K}uramoto {MFG}
  with strong interaction.
\newblock \emph{Communications in Partial Differential Equations}
  49 (2024), no.\ 1--2, 121--147. arXiv:2307.09305.

\bibitem[CG20]{CG}
R.~Carmona and C.~V.~Graves.
\newblock Jet lag recovery: synchronization of circadian oscillators
  as a mean field game.
\newblock \emph{Dynamic Games and Applications} 10 (2020), no.\ 1,
  79--99.

\bibitem[Cir19]{Cirant}
M.~Cirant.
\newblock On the existence of oscillating solutions in non-monotone
  mean-field games.
\newblock \emph{Journal of Differential Equations} 266 (2019), no.\ 12,
  8067--8093.

\bibitem[GPP12]{GPP}
G.~Giacomin, K.~Pakdaman, and X.~Pellegrin.
\newblock Global attractor and asymptotic dynamics in the {K}uramoto
  model for coupled noisy phase oscillators.
\newblock \emph{Nonlinearity} 25 (2012), no.\ 5, 1247--1273.

\bibitem[GT01]{GT}
D.~Gilbarg and N.~S.~Trudinger.
\newblock \emph{Elliptic Partial Differential Equations of Second Order}.
\newblock Reprint of the 1998 edition. Classics in Mathematics.
Springer, Berlin, 2001.

\bibitem[HMC06]{HMC}
M.~Huang, R.~P.~Malham\'e, and P.~E.~Caines.
\newblock Large population stochastic dynamic games: closed-loop
  {M}c{K}ean--{V}lasov systems and the {N}ash certainty equivalence
  principle.
\newblock \emph{Communications in Information \& Systems} 6 (2006),
  no.\ 3, 221--252.

\bibitem[HS25]{HS}
F.~H\"ofer and H.~M.~Soner.
\newblock Synchronization games.
\newblock \emph{Mathematics of Operations Research} 51 (2025),
  no.\ 2, 1443--1462. arXiv:2402.08842.

\bibitem[Kur75]{Kuramoto}
Y.~Kuramoto.
\newblock Self-entrainment of a population of coupled non-linear
  oscillators.
\newblock In \emph{International Symposium on Mathematical Problems in
  Theoretical Physics}, Lecture Notes in Physics, vol.\ 39, pages
  420--422. Springer, 1975.

\bibitem[LL07]{LL}
J.-M.~Lasry and P.-L.~Lions.
\newblock Mean field games.
\newblock \emph{Japanese Journal of Mathematics} 2 (2007), no.\ 1,
  229--260.

\bibitem[Str00]{Strogatz}
S.~H.~Strogatz.
\newblock From {K}uramoto to {C}rawford: exploring the onset of
  synchronization in populations of coupled oscillators.
\newblock \emph{Physica D: Nonlinear Phenomena} 143 (2000), no.\ 1--4,
  1--20.

\bibitem[YMMS12]{YMMS}
H.~Yin, P.~G.~Mehta, S.~P.~Meyn, and U.~V.~Shanbhag.
\newblock Synchronization of coupled oscillators is a game.
\newblock \emph{IEEE Transactions on Automatic Control} 57 (2012),
  no.\ 4, 920--935.

\end{thebibliography}

\end{document}
